% appC2 —— 附录 C 后段：Infinite-dimensional case（原文 wcf-paper.tex 3686–3822 行）
% 原文小节标题为 \subsection*{Infinite-dimensional case}（不编号；TOC 中亦无 C.1/C.2 条目），
% 故此处保持带星号，与原文排版及上一块的 \subsection*{有限维情形} 一致。
% 原文宏 \eps 不在本工程 preamble 中，此处本地补齐（与原 wcf-paper.tex 的定义相同）。
\providecommand{\eps}{\epsilon}

\subsection*{无穷维情形}

现在我们对乘性积分方程 式 (5.13) 给出类似的分析。% TODO 合稿改 \eqref{eq:X-integral-mult-explicit}
首先，对任意向量 $\gamma \in \clat$，我们用如下幂级数展开来定义一个向量 $f^\gamma \in \clat_\IQ$：
\begin{equation}
- \sum_{\gamma' \in \clat} \left( \degen(\gamma';u)
\log (1- \sigma(\gamma')
\CX_{\gamma'} ) \right) \gamma' = \sum_{\gamma' \in \clat} f^{\gamma'} \CX_{\gamma'}
\end{equation}
或者更显式地，
\begin{equation}
f^\gamma = \sum_{n \ge 1 \text{ s.t. } \gamma = n \gamma'} \frac{\sigma(\gamma')^n}{n} \degen(\gamma';u) \gamma'.
\end{equation}
这样定义的用意在于：此后 式 (5.13) 便取如下形式% TODO 合稿改 \eqref{eq:X-integral-mult-explicit}
\begin{equation} \label{eq:X-integral-mult-std}
\CX_\gamma(\zeta) = \CX^{\sfm}_\gamma(\zeta) \exp \left\langle \gamma,
\frac{1}{4 \pi i} \sum_{\gamma'} f^{\gamma'} \int_{\ell_{\gamma'}} \frac{d\zeta'}{\zeta'} \frac{\zeta' +
\zeta}{\zeta' - \zeta} \CX_{\gamma'}(\zeta') \right\rangle.
\end{equation}

我们希望把 式 \eqref{eq:X-integral-mult-std} 的解 $\CX$ 构造为逐次逼近 $\CX^{(n)}$ 的极限，或者构造为 式 (5.1) 中所定义的 $\Upsilon$ 的相应逼近 $\Upsilon^{(n)}$ 的极限。% TODO 合稿改 \eqref{eq:def-upsilon}
$\Upsilon$ 是从 $\CM_u$ 到其复化的一个映射；我们把它的分量写成函数，即 $\Upsilon_\gamma(\theta) := \gamma \cdot \Upsilon(\theta)$。
我们取 $\Upsilon^{(0)}_\gamma = \theta_\gamma$ 作为起点。
回忆 $\Upsilon$ 的定义方式使下式成立
\begin{equation}
\CX^{\sfm}(\Upsilon^{(n)}) = \CX^{(n)}(\theta),
\end{equation}
于是迭代步可写成
\begin{equation} \label{eq:recursion}
e^{i \Upsilon^{(n+1)}_\gamma} = e^{i \theta_\gamma} \exp \left\langle \gamma,
\frac{1}{4 \pi i} \sum_{\gamma'} f^{\gamma'} \int_{\ell_{\gamma'}}
\frac{d\zeta'}{\zeta'} \frac{\zeta' + \zeta}{\zeta' - \zeta}
\CX^{\sfm}_{\gamma'}(\Upsilon^{(n)},\zeta')\right\rangle
\end{equation}
该迭代的不动点 $\Upsilon^{(n+1)} = \Upsilon^{(n)}$ 将给出 式 (5.11) 的一个解。% TODO 合稿改 \eqref{eq:X-integral-mult}
因此，要看出解确实存在，我们应当验证这一迭代是压缩的，即
\begin{equation}
\max_{\zeta,\theta} \norm{\Upsilon^{(n+1)} - \Upsilon^{(n)}} < C \max_{\zeta,\theta} \norm{\Upsilon^{(n)} - \Upsilon^{(n-1)}}
\end{equation}
其中 $C < 1$ 为某常数。
更确切地说，我们将验证：当这一迭代作用在满足 $\max_{\zeta,\theta} \norm{\Upsilon^{(n)} - \theta} < \infty$、并且还满足一个附加条件的诸 $\Upsilon^{(n)}$ 上时，它是压缩的；该附加条件表达的正是「它们离实环面不太远」这一事实：我们要求
$\max_{\zeta,\theta} |e^{i \Upsilon^{(n)}_\gamma}|<e^{\eps \norm{\gamma}}$，其中常数 $\eps>0$ 稍后确定。

首先我们要看出，这一迭代保持我们的附加条件。
取 式 \eqref{eq:recursion} 的绝对值并作鞍点分析，当 $R$ 足够大时便得到估计\footnote{与上一小节一样，当 $\zeta$ 恰好落在鞍点上时，这一分析还须补上一段单独的讨论；但那只会使压低程度相差一个 $\sqrt{R}$，仍足以让我们建立 式 \eqref{eq:estimate}。}
\begin{equation} \label{eq:estimate}
|e^{i \Upsilon^{(n+1)}_\gamma}| <  \exp \left[ \sum_{\gamma'} \abs{\inprod{f^{\gamma'},\gamma}}
e^{- 2 \pi R |Z_{\gamma'}| + \eps \norm{\gamma'}} \right].
\end{equation}
现在假设 $\sum_{\gamma'} f^{\gamma'} \norm{\gamma'} e^{- 2 \pi R |Z_{\gamma'}|}$ 在 $R$ 足够大时收敛。（这相当于要求诸 $\degen(\gamma';u)$ 随 $\gamma'$ 增长得不太快；对场论而言这看来非常合理，但在引力情形几乎肯定会被破坏。）
我们还利用第 2.2 节所回顾的支撑性质（Support Property），用 $K \abs{Z_{\gamma'}}$ 给出 $\norm{\gamma'}$ 的上界。% TODO 合稿改 \ref{sec:ReviewKS}
于是当 $R$ 足够大时，我们可以选取 $\eps$，使右端小于 $e^{\eps \norm{\gamma}}$。
对 $\gamma$ 取遍 $\clat$ 的一组基逐一这样做，就得到所需的 $\eps$。（事实上，当 $R$ 很大时，可以取 $\epsilon \to 0$ 以指数速度变快。）

现在我们想估计
\begin{align}
\norm{\Upsilon^{(n+1)}-\Upsilon^{(n)}} &=
\frac{1}{4 \pi} \bigg\lVert \sum_{\gamma' \in \clat} \int_{\ell_{\gamma'}}
\frac{d\zeta'}{\zeta'} \frac{\zeta' + \zeta}{\zeta' - \zeta}
f^{\gamma'}\left(\CX^{\sfm}_{\gamma'}(\Upsilon^{(n)},\zeta')-\CX^{\sfm}_{\gamma'}(\Upsilon^{(n-1)},\zeta')\right) \bigg\rVert \\
&\le \frac{1}{4 \pi}\bigg\lVert \sum_{\gamma' \in \clat} \int_{\ell_{\gamma'}}
\frac{d\zeta'}{\zeta'} \frac{\zeta' + \zeta}{\zeta' - \zeta}
f^{\gamma'} \abs{\CX^{\sfm}_{\gamma'}(\theta, \zeta')} \left(e^{i \Upsilon^{(n)}_{
\gamma'}(\zeta')} - e^{i \Upsilon^{(n-1)}_{\gamma'}(\zeta')}\right) \bigg\rVert \\
&\le \frac{1}{4 \pi}\bigg\lVert \sum_{\gamma' \in \clat} \int_{\ell_{\gamma'}}
\frac{d\zeta'}{\zeta'} \frac{\zeta' + \zeta}{\zeta' - \zeta}
f^{\gamma'} \abs{\CX^{\sfm}_{\gamma'}(\theta, \zeta')}  \bigg\rVert e^{\eps \norm{\gamma'}} \norm{\gamma'} \max_{\zeta}  \norm{\Upsilon^{(n)} -
\Upsilon^{(n-1)}}.
\end{align}
随后，大 $R$ 下的鞍点分析给出
\begin{equation}
\norm{\Upsilon^{(n+1)}-\Upsilon^{(n)}} \le
\frac{1}{4\pi}\max_{\zeta} \norm{\Upsilon^{(n)}(\zeta) -
\Upsilon^{(n-1)}(\zeta)} \bigg\lVert \sum_{\gamma' \in \clat}
f^{\gamma'} \norm{\gamma'} e^{- 2 \pi R |Z_{\gamma'}| + \eps
\norm{\gamma'}} \bigg\rVert.
\end{equation}
当 $R$ 足够大时，在我们的收敛假设下，这就建立了压缩性质；事实上该迭代收敛得很快，其速度由最大的那个 $e^{- 2 \pi R |Z_{\gamma'}|}$ 决定。

我们可以用与有限维情形中 $\CF$ 同类的函数，写出 $\CX_\gamma(\zeta)$ 的显式表达式。
递归关系中指数因子的出现使事情稍微复杂一些：现在不再对链求和，而是对带装饰的有根树求和。
设 $\CT$ 是一棵有根树，其边用数对 $(i,j)$ 标记（其中 $i$ 是更靠近根节点的那个节点），每个节点则带一个装饰，即取定 $\clat$ 中的一个 $\gamma_i$。
此外，把根节点处的装饰记为 $\gamma_\CT$。
于是定义该树的权重为 $\clat_\IQ$ 中的一个元素，它由诸 $\degen(\gamma;u)$ 确定：
\begin{equation}
\CW_\CT = \frac{f^{\gamma_\CT}}{\abs{\mathrm{Aut}(\CT)}} \prod_{(i,j) \in {\mathrm {Edges}}(\CT)} \inprod{\gamma_i, f^{\gamma_j}}.
\end{equation}
于是 $\CX_\gamma(\zeta)$ 的迭代解取如下形式
\begin{equation} \label{eq:iterative-X}
\CX_\gamma(\zeta) = \CX^{\sfm}_\gamma (\zeta)\exp \left\langle \gamma, \sum_\CT
\CW_\CT \CG_\CT(\zeta) \right\rangle,
\end{equation}
其中 $\CG_\CT(\zeta)$ 为某些函数。关于 $\CX$ 的积分方程 式 \eqref{eq:X-integral-mult-std} 则变成一条公式，把每个 $\CG_\CT(\zeta)$ 用较小的树所对应的那些函数表示出来。
具体说来，从 $\CT$ 中删去根节点，会留下一组有根树 $\CT_a$；而 式 \eqref{eq:X-integral-mult-std} 将被满足，只要
\begin{equation}
\CG_\CT(\zeta) = \frac{1}{4 \pi i} \int_{\ell_{\gamma_\CT}} \frac{d \zeta'}{\zeta'} \frac{\zeta' + \zeta}{\zeta' - \zeta} \CX^{\sfm}_{\gamma_\CT}(\zeta') \prod_a \CG_{\CT_a}(\zeta').
\end{equation}
由此可见，与有限维情形中的 $\CF^{(n)}$ 一样，$\CG_\CT$ 沿 $\ell_{\gamma_\CT}$ 的不连续性由较低的那些 $\CG_{\CT_a}$ 之乘积决定。

由 式 \eqref{eq:iterative-X} 我们还能直接得到 $\Upsilon$：
\begin{equation}
e^{i \Upsilon_\gamma} = \exp \left\langle \gamma, \sum_\CT
\CW_\CT \CG_\CT(\zeta) \right\rangle.
\end{equation}
特别地，这使我们能够计算 $\Upsilon(\zeta = 0)$。
辛形式 $\hsymp(\zeta)$ 的展开也可用同样的方式分析，并组织成对树的求和。
